\documentclass[12pt]{article}
\usepackage{amsmath, amssymb}
\usepackage{geometry}
\usepackage{tcolorbox}
\usepackage{tikz}
\usepackage{pgfplots}
\pgfplotsset{compat=1.18}
\usepackage{multicol}
\usepackage{enumitem}
\usepackage{xcolor}
\usepackage{booktabs}
\usepackage{array}

\geometry{margin=1in}
\tcbuselibrary{skins, breakable}

%-------------------------------------------------
% Colors & box styles
%-------------------------------------------------
\definecolor{mainblue}{RGB}{25, 70, 160}
\definecolor{lightblue}{RGB}{218, 232, 255}
\definecolor{lightgray}{RGB}{246, 246, 246}
\definecolor{darkgray}{RGB}{70, 70, 70}
\definecolor{alertred}{RGB}{170, 20, 20}
\definecolor{green1}{RGB}{15, 110, 45}
\definecolor{orange1}{RGB}{200, 90, 0}

\tcbset{
  ansbox/.style={
    colback=lightgray, colframe=darkgray!60,
    fonttitle=\bfseries, title={Answer Key},
    rounded corners, boxrule=0.8pt, breakable
  }
}

%-------------------------------------------------
% Helpers
%-------------------------------------------------
\newcommand{\blank}{\underline{\hspace{2.5cm}}}

%-------------------------------------------------
\title{%
  \textbf{\Large Chapter 11 Mid-Chapter Quiz}\\[4pt]
  \large Quadratic Functions\\[2pt] 
}
\author{FangMath}
\date{Name: \underline{\hspace{6cm}} \quad Date: \underline{\hspace{3.5cm}}}

\begin{document}
\maketitle
\vspace{-6pt}
\noindent\rule{\linewidth}{0.6pt}

%=======================================================
\section*{\textcolor{green1}{Level 1 — Basic Concepts}}
%=======================================================

%-----------
\subsection*{Part A: Identify the Vertex and Axis of Symmetry}
%-----------

\noindent Graph each function. Label the \textbf{axis of symmetry} and the \textbf{vertex}.

\begin{enumerate}[leftmargin=*, label=\textbf{\arabic*.}]

  \item $y = \dfrac{1}{2}x^2$
  \vspace{3.5cm}

  \item $y = -2x^2 - 1$
  \vspace{3.5cm}

  \item $y = 3x^2 - 6x$
  \vspace{3.5cm}

  \item $y = x^2 + 2x + 4$
  \vspace{3.5cm}

  \item $y = -0.5x^2 + 2x + 1$
  \vspace{3.5cm}

\end{enumerate}

%-----------
\subsection*{Part B: Order by Width of Graph}
%-----------

\noindent Order each group of quadratic functions from \textbf{widest} to \textbf{narrowest} graph.

\begin{enumerate}[leftmargin=*, label=\textbf{\arabic*.}, start=6]

  \item $y = 2x^2, \quad y = 0.5x^2, \quad y = -x^2$

  \vspace{1cm}
  Answer: \blank

  \item $f(x) = 4x^2, \quad f(x) = \dfrac{2}{3}x^2, \quad f(x) = 3x^2$

  \vspace{1cm}
  Answer: \blank

  \item $f(x) = 0.6x^2, \quad f(x) = 0.3x^2, \quad f(x) = 0.2x^2$

  \vspace{1cm}
  Answer: \blank

  \item $y = -2x^2, \quad y = x^2, \quad y = -0.25x^2$

  \vspace{1cm}
  Answer: \blank

\end{enumerate}

%=======================================================
\section*{\textcolor{orange1}{Level 2 — Solving Quadratics}}
%=======================================================

%-----------
\subsection*{Part C: Solve by Finding Square Roots}
%-----------

\noindent Solve each equation by finding square roots. If the equation has \textbf{no real-number solution}, write \textit{no solution}.

\begin{enumerate}[leftmargin=*, label=\textbf{\arabic*.}, start=10]

  \item $m^2 = 81$

  \vspace{2.2cm}

  \item $5r^2 - 180 = 0$

  \vspace{2.2cm}

  \item $t^2 - 7 = -18$

  \vspace{2.2cm}

  \item $36n^2 = 9$

  \vspace{2.2cm}

\end{enumerate}

%-----------
\subsection*{Part D: Solve by Graphing}
%-----------

\noindent Solve each equation by \textbf{graphing the related function}. If the equation has no real-number solution, write \textit{no solution}.

\begin{enumerate}[leftmargin=*, label=\textbf{\arabic*.}, start=14]

  \item $x^2 - 16 = 0$
  \vspace{3.5cm}

  \item $x^2 + 9 = 0$
  \vspace{3.5cm}

  \item $0.25x^2 = 0$
  \vspace{3.5cm}

\end{enumerate}

%-----------
\subsection*{Part E: Solve by Factoring}
%-----------

\noindent Solve each equation by \textbf{factoring}.

\begin{enumerate}[leftmargin=*, label=\textbf{\arabic*.}, start=17]

  \item $b^2 + 3b - 4 = 0$

  \vspace{2.5cm}

  \item $n^2 + n - 12 = 0$

  \vspace{2.5cm}

  \item $2x^2 - 5x - 3 = 0$

  \vspace{2.5cm}

  \item $t^2 - 3t = 28$

  \vspace{2.5cm}

  \item $3n^2 = 6n$

  \vspace{2.5cm}

\end{enumerate}

%=======================================================
\section*{\textcolor{alertred}{Level 3 — Advanced \& Word Problems}}
%=======================================================

%-----------
\subsection*{Part F: Conceptual Questions}
%-----------

\begin{enumerate}[leftmargin=*, label=\textbf{\arabic*.}, start=22]

  \item \textbf{Writing.} Describe the steps you would use to graph the function $y = 2x^2 + 5$.
  \vspace{3.5cm}

  \item \textbf{Reasoning.} Does the value of $c$ in the quadratic function $y = ax^2 + bx + c$
        affect the \emph{horizontal position} of the vertex of the graph? Explain why or why not.
  \vspace{3.5cm}

  \item \textbf{Writing.} Describe how the graph of $y = 3x^2$ differs from the graph of $y = x^2$.
  \vspace{3cm}

\end{enumerate}

%-----------
\subsection*{Part G: Open-Ended}
%-----------

\noindent Give an example of a quadratic function that matches each description.

\begin{enumerate}[leftmargin=*, label=\textbf{\arabic*.}, start=25]

  \item The axis of symmetry is to the \textbf{left} of the $y$-axis.

  \vspace{1.5cm}
  $f(x) = $ \blank

  \item Its graph lies \textbf{entirely below} the $x$-axis.

  \vspace{1.5cm}
  $f(x) = $ \blank

  \item Its graph opens \textbf{upward} and has its vertex at $(0,0)$.

  \vspace{1.5cm}
  $f(x) = $ \blank

\end{enumerate}

%-----------
\subsection*{Part H: Extended Problem}
%-----------

\begin{enumerate}[leftmargin=*, label=\textbf{\arabic*.}, start=28]

  \item \textbf{a.} Solve $x^2 - 4 = 0$ and $2x^2 - 8 = 0$ by graphing their related functions.

  \vspace{4cm}

  \textbf{b.} \textbf{Reasoning.} Why does it make sense that the graphs have the same $x$-intercepts?

  \vspace{3cm}

  \item \textbf{(STEM — Construction)} You are building a rectangular planter for your school
        garden. You want the area of the bottom to be $90\ \text{ft}^2$. You want the length of
        the planter to be 3 ft longer than \emph{twice} its width. What should the dimensions of
        the bottom of the planter be?

  \vspace{5cm}

  \item \textbf{(Sewing)} You have $324\ \text{ft}^2$ of fabric to make a circular play parachute
        for kids. What is the radius of the largest parachute you could make? Round to the nearest
        tenth of a foot. \textit{(Hint: $A = \pi r^2$)}

  \vspace{4cm}

\end{enumerate}

%=======================================================
\newpage
\section*{Answer Key}
%=======================================================

\begin{tcolorbox}[ansbox]

\textbf{Part A — Graphs (vertex and axis of symmetry)}

\begin{enumerate}[leftmargin=*, label=\textbf{\arabic*.}]
  \item $y = \frac{1}{2}x^2$: vertex $(0,0)$, axis $x = 0$. Opens upward, wider than $y=x^2$.
  \item $y = -2x^2 - 1$: vertex $(0,-1)$, axis $x = 0$. Opens downward.
  \item $y = 3x^2 - 6x = 3(x-1)^2 - 3$: vertex $(1,-3)$, axis $x = 1$.
  \item $y = x^2 + 2x + 4 = (x+1)^2 + 3$: vertex $(-1,3)$, axis $x = -1$.
  \item $y = -0.5x^2 + 2x + 1 = -0.5(x-2)^2 + 3$: vertex $(2,3)$, axis $x = 2$.
\end{enumerate}

\medskip
\textbf{Part B — Widest to Narrowest} (smaller $|a|$ = wider graph)

\begin{enumerate}[leftmargin=*, label=\textbf{\arabic*.}, start=6]
  \item $y=0.5x^2,\; y=-x^2,\; y=2x^2$
  \item $f(x)=\frac{2}{3}x^2,\; f(x)=3x^2,\; f(x)=4x^2$
  \item $f(x)=0.2x^2,\; f(x)=0.3x^2,\; f(x)=0.6x^2$
  \item $y=-0.25x^2,\; y=x^2,\; y=-2x^2$
\end{enumerate}

\medskip
\textbf{Part C — Square Roots}

\begin{enumerate}[leftmargin=*, label=\textbf{\arabic*.}, start=10]
  \item $m^2 = 81 \Rightarrow m = \pm 9$
  \item $5r^2 = 180 \Rightarrow r^2 = 36 \Rightarrow r = \pm 6$
  \item $t^2 = -11$: \textit{no solution} (negative under square root)
  \item $36n^2 = 9 \Rightarrow n^2 = \frac{1}{4} \Rightarrow n = \pm\frac{1}{2}$
\end{enumerate}

\medskip
\textbf{Part D — Graphing}

\begin{enumerate}[leftmargin=*, label=\textbf{\arabic*.}, start=14]
  \item $x^2 = 16 \Rightarrow x = \pm 4$
  \item $x^2 = -9$: \textit{no solution}
  \item $0.25x^2 = 0 \Rightarrow x = 0$
\end{enumerate}

\medskip
\textbf{Part E — Factoring}

\begin{enumerate}[leftmargin=*, label=\textbf{\arabic*.}, start=17]
  \item $(b+4)(b-1)=0 \Rightarrow b = -4,\; 1$
  \item $(n+4)(n-3)=0 \Rightarrow n = -4,\; 3$
  \item $(2x+1)(x-3)=0 \Rightarrow x = -\frac{1}{2},\; 3$
  \item $t^2-3t-28=0 \Rightarrow (t-7)(t+4)=0 \Rightarrow t = 7,\; -4$
  \item $3n^2-6n=0 \Rightarrow 3n(n-2)=0 \Rightarrow n = 0,\; 2$
\end{enumerate}

\medskip
\textbf{Part F — Conceptual}

\begin{enumerate}[leftmargin=*, label=\textbf{\arabic*.}, start=22]
  \item Graph $y = 2x^2 + 5$: vertex at $(0,5)$; axis of symmetry $x=0$; opens upward; narrower than $y=x^2$ (since $a=2>1$); $y$-intercept at $(0,5)$.
  \item $c$ does \textbf{not} affect the horizontal position of the vertex. The $x$-coordinate of the vertex is $x = -\frac{b}{2a}$, which depends only on $a$ and $b$. Changing $c$ shifts the graph vertically (up or down) but not horizontally.
  \item $y = 3x^2$ is \textbf{narrower} than $y = x^2$ because $|a| = 3 > 1$. Both open upward with vertex $(0,0)$, but $y=3x^2$ rises more steeply.
\end{enumerate}

\medskip
\textbf{Part G — Open-Ended} (sample answers)

\begin{enumerate}[leftmargin=*, label=\textbf{\arabic*.}, start=25]
  \item e.g.\ $f(x) = (x+3)^2$, axis $x = -3 < 0$
  \item e.g.\ $f(x) = -x^2 - 1$ (always $\le -1 < 0$)
  \item e.g.\ $f(x) = 5x^2$ (opens up, vertex at origin)
\end{enumerate}

\medskip
\textbf{Part H — Extended}

\begin{enumerate}[leftmargin=*, label=\textbf{\arabic*.}, start=28]
  \item \textbf{a.} Both equations give $x = \pm 2$.\\
        \textbf{b.} $2x^2 - 8 = 2(x^2-4)$. Multiplying by a constant scales the graph vertically but does not change the $x$-intercepts (where $y=0$).

  \item Let width $= w$. Then length $= 2w+3$.
        \[
          w(2w+3) = 90 \implies 2w^2+3w-90=0
        \]
        \[
          w = \frac{-3 + \sqrt{9+720}}{4} = \frac{-3+27}{4} = 6\ \text{ft}, \quad
          \ell = 2(6)+3 = 15\ \text{ft}.
        \]
        \textbf{Dimensions: $6\ \text{ft} \times 15\ \text{ft}$.}

  \item $\pi r^2 = 324 \Rightarrow r^2 = \dfrac{324}{\pi} \approx 103.13 \Rightarrow r \approx \mathbf{10.2\ \text{ft}}$.
\end{enumerate}

\end{tcolorbox}

\vfill
\begin{center}
\textit{\textcopyright\ FangMath --- All rights reserved --- 2026}
\end{center}

\end{document}
